\documentclass{article}
\usepackage[utf8]{inputenc}
\usepackage{amsmath,amssymb}
\usepackage[most]{tcolorbox}
\usepackage{geometry}
\geometry{margin=2.5cm}

\newcommand{\step}[1]{\par\noindent\hangindent=3.5em\hangafter=1%
  \makebox[3.5em][l]{\textbf{#1.}}}
\newcommand{\steptag}[1]{\hfill{\small\textsf{[#1]}}}

\newtcolorbox{proofbox}[1]{
  colback=gray!5, colframe=gray!60,
  fonttitle=\bfseries, title={#1},
  breakable, enhanced,
  left=4pt, right=4pt, top=4pt, bottom=4pt,
  before skip=10pt, after skip=10pt
}

\begin{document}

\begin{proofbox}{There are universal K\_step >= 1 and e\_step > 0 such that,...}
\step{1} There are universal K\_step >= 1 and e\_step > 0 such that, if B is a finite-dimensional C*-algebra, A an extended epsilon-C*-algebra, and v:B->A an extended d-inclusion with d+epsilon <= e\_step, then one dagger-preserving level-one map v\textasciicircum{}+, with v\_n\textasciicircum{}+ = I\_n tensor v\textasciicircum{}+, satisfies sup\_n ||v\_n\textasciicircum{}+ - v\_n|| <= K\_step*d and is an extended d\textasciicircum{}+-inclusion for d\textasciicircum{}+ <= K\_step*(d\textasciicircum{}2+epsilon). \steptag{validated} \\[6pt]
\step{1.1} Every finite-dimensional C*-algebra B has a diagonal D=sum\_j A\_j tensor B\_j in the sense of def-fd-cstar-diagonal with sum\_j ||A\_j|| ||B\_j||=1; this representation is dimension-free. \steptag{validated} \\[6pt]
\step{1.1.1} Let C be a finite-dimensional unital C*-algebra. First note the following elementary finite-dimensional spectral fact, used below rather than any classification theorem. If a=a\textasciicircum{}dagger, then exp(ita) is unitary, so spectral mapping forces spectrum(a) to be real; the spectral-radius formula together with the C*-identity gives ||a||=r(a). Hence the polynomial product over the distinct spectral values annihilates a, and Lagrange interpolation gives a=sum\_lambda lambda p\_lambda with pairwise orthogonal self-adjoint projections p\_lambda summing to 1. In particular spectral support projections are polynomials with zero constant term.

For every nonzero dagger-ideal I of C choose a basis x\_1,...,x\_N and put h=sum\_k(x\_k\textasciicircum{}dagger x\_k+x\_k x\_k\textasciicircum{}dagger). Let z be the spectral support projection of h. Since z is a polynomial in h with zero constant term, z lies in I. From h(1-z)=0 and positivity of each summand, (1-z)x\_k\textasciicircum{}dagger x\_k(1-z)=0 and (1-z)x\_k x\_k\textasciicircum{}dagger(1-z)=0; the C*-identity gives x\_k(1-z)=0=(1-z)x\_k. Thus z is the two-sided identity of I. For c in C, both cz and zc lie in I, whence zcz=cz and zcz=zc, so z is central.

Apply the spectral fact inside Z(C) and repeatedly take a nonzero projection minimal among central projections. This gives orthogonal minimal central projections z\_1,...,z\_m with sum z\_r=1 and C=direct\_sum\_r C\_r, C\_r=z\_r C. Any central self-adjoint element of C\_r has central spectral projections below z\_r and is therefore scalar; hence Z(C\_r)=C z\_r. If J is a nonzero dagger-ideal of C\_r, the preceding ideal argument gives it a nonzero central identity, necessarily z\_r, so J=C\_r. Thus each C\_r is simple.

Fix one simple C\_r, writing its unit as 1. A nonzero projection exists by taking the spectral support of y\textasciicircum{}dagger y for nonzero y. Choose a nonzero projection p minimizing dim(p C\_r p). If p C\_r p were not C p, a non-scalar self-adjoint element in it would have a nonzero proper spectral projection q<p and q C\_r q would have smaller dimension. Hence p C\_r p=C p. Since C\_r p C\_r is a nonzero dagger-ideal, simplicity gives C\_r p C\_r=C\_r. Set v\_1=p and p\_1=p. Inductively, suppose v\_i\textasciicircum{}dagger v\_i=p and the projections p\_i=v\_i v\_i\textasciicircum{}dagger are pairwise orthogonal. If e=1-sum\_i p\_i is nonzero, an expression 1=sum\_l a\_l p b\_l implies some x=e a\_l p is nonzero. Then x\textasciicircum{}dagger x belongs to p C\_r p, so x\textasciicircum{}dagger x=lambda p with lambda>0. For v=lambda\textasciicircum{}(-1/2)x one has v\textasciicircum{}dagger v=p and vv\textasciicircum{}dagger is a projection below e. Adjoining v therefore extends the family. Finite dimensionality makes the process terminate with sum\_i p\_i=1.

Put e\_ij=v\_i v\_j\textasciicircum{}dagger. Orthogonality gives v\_j\textasciicircum{}dagger v\_k=delta\_jk p, hence e\_ij e\_kl=delta\_jk e\_il and e\_ij\textasciicircum{}dagger=e\_ji. Moreover, for a in C\_r, v\_i\textasciicircum{}dagger a v\_j lies in p C\_r p=Cp, so p\_i a p\_j is a scalar multiple of e\_ij. As a=sum\_ij p\_i a p\_j, the e\_ij span C\_r. Therefore E\_ij maps to e\_ij is a bijective unital dagger-homomorphism M\_q to C\_r. It preserves spectra algebraically; applying ||b||=r(b) to the positive elements b=a\textasciicircum{}dagger a and using the C*-identity proves ||Phi(a)||\textasciicircum{}2=r(Phi(a\textasciicircum{}dagger a))=r(a\textasciicircum{}dagger a)=||a||\textasciicircum{}2. Thus the identification, and consequently the central direct-sum identification, is *-isometric.

It remains to construct the diagonal in a block without further structure input. In M\_q index matrix units modulo q, let omega=exp(2 pi i/q), S xi\_k=xi\_(k+1), T xi\_k=omega\textasciicircum{}k xi\_k, and W\_ab=S\textasciicircum{}a T\textasciicircum{}b. These are unitaries and character orthogonality gives
q\textasciicircum{}(-2) sum\_(a,b) W\_ab\textasciicircum{}dagger tensor W\_ab = q\textasciicircum{}(-1) sum\_(i,j) E\_ij tensor E\_ji=:D\_q.
For X=E\_rs, direct multiplication gives X D\_q=q\textasciicircum{}(-1)sum\_j E\_rj tensor E\_js=D\_q X, hence the equality for every X; also pi(D\_q)=q\textasciicircum{}(-1)sum\_(i,j)E\_ij E\_ji=I. The displayed Weyl expression has q\textasciicircum{}2 positive coefficients q\textasciicircum{}(-2), each multiplying U\textasciicircum{}dagger tensor U with U unitary, so the sum of the corresponding norm products is exactly 1. \steptag{validated} \\[4pt]
\step{1.1.2} For B=direct\_sum\_(r=1)\textasciicircum{}m M\_(q\_r), independently choose in each block a Pauli unitary W\_(r,a,b) and a sign sigma\_r in \{+1,-1\}, and average U\textasciicircum{}dagger tensor U over U=direct\_sum\_r sigma\_r W\_(r,a,b). Independent sign averaging kills every cross-block tensor component, while within each block it leaves D\_(q\_r); hence the average is a diagonal of B. It is a convex combination of U\textasciicircum{}dagger tensor U with U unitary, so its displayed representation has sum\_j||A\_j||||B\_j||=1, independent of m and all q\_r. \steptag{validated} \\[4pt]
\step{1.2} Fix such a diagonal. For g(X,Y)=v(XY)-v(X)v(Y), define w\_prime(X)=sum\_j v(A\_j)g(B\_j,X), w\_doubleprime(X)=w\_prime(X\textasciicircum{}dagger)\textasciicircum{}dagger, w=(w\_prime+w\_doubleprime)/2, and v\textasciicircum{}+=v+w. For all n, I\_n tensor w is exactly the correction obtained from v\_n and I\_n tensor D, and there is a universal C\_w such that sup\_n ||I\_n tensor w|| <= C\_w d whenever d+epsilon is sufficiently small. \steptag{validated} \\[6pt]
\step{1.2.1} For X=[X\_pq] in M\_n tensor B, the defining matrix product gives [g\_n(I\_n tensor B\_j,X)]\_pq=g(B\_j,X\_pq), and then [sum\_j v\_n(I\_n tensor A\_j)g\_n(I\_n tensor B\_j,X)]\_pq=w\_prime(X\_pq). The adjoint rule gives the same identity for w\_doubleprime. Therefore the correction constructed at every level is exactly I\_n tensor w and v\_n\textasciicircum{}+=I\_n tensor v\textasciicircum{}+; no level-dependent choice is made. \steptag{validated} \\[4pt]
\step{1.2.2} Invoke the registered external GT-kitaev-standard-diagonal: choose D=sum\_j A\_j tensor B\_j with A\_j=p\_j U\_j\textasciicircum{}dagger, B\_j=U\_j, p\_j>=0, sum\_j p\_j=1, and U\_j unitary. Thus D is a diagonal in the sense of def-fd-cstar-diagonal and sum\_j||A\_j||||B\_j||=1, with no classification theorem or sibling result being used. For X=[X\_pq] in M\_n tensor B, define g\_n(S,T)=v\_n(ST)-v\_n(S)v\_n(T) and W\_prime\_n(X)=sum\_j v\_n(I\_n tensor A\_j)g\_n(I\_n tensor B\_j,X). Direct matrix multiplication gives [g\_n(I\_n tensor B\_j,X)]\_pq=g(B\_j,X\_pq), hence [W\_prime\_n(X)]\_pq=w\_prime(X\_pq), so W\_prime\_n=I\_n tensor w\_prime; applying the level-n involution gives W\_doubleprime\_n=I\_n tensor w\_doubleprime. By the block-diagonal axiom in def-operator-space, ||I\_n tensor A\_j||=||A\_j|| and ||I\_n tensor B\_j||=||B\_j||. Unpacking def-extended-delta-inclusion at level n gives ||v\_n(S)||<=(1+d)||S|| and, through its d-homomorphism clause, ||g\_n(S,T)||<=d||S||||T||. Unpacking def-extended-epsilon-cstar-algebra through the registered def-epsilon-cstar-algebra gives ||ST||<=(1+epsilon)||S||||T|| and an isometric involution. Consequently ||W\_prime\_n(X)|| <= (1+epsilon)(1+d)d sum\_j||A\_j||||B\_j||||X||=(1+epsilon)(1+d)d||X||, and the same bound holds for W\_doubleprime\_n. Since I\_n tensor w=(W\_prime\_n+W\_doubleprime\_n)/2 and d+epsilon<=1/4 implies (1+epsilon)(1+d)<=25/16<2, one has ||I\_n tensor w||<=2d uniformly in n and B. \steptag{validated} \\[4pt]
\step{1.3} The same correction satisfies, at every amplification, ||v\_n\textasciicircum{}+(XY)-v\_n\textasciicircum{}+(X)v\_n\textasciicircum{}+(Y)|| <= C\_m(d\textasciicircum{}2+epsilon)||X||||Y|| for a universal C\_m independent of B, A, n, d, and epsilon. \steptag{validated} \\[6pt]
\step{1.3.1} Fix n. The following inequalities are not extra hypotheses: def-extended-epsilon-cstar-algebra says M\_n tensor A is an epsilon-C*-algebra, and registered def-epsilon-cstar-algebra therefore gives ||ST||<=(1+epsilon)||S||||T||, ||(ST)R-S(TR)||<=epsilon||S||||T||||R||, the isometric anti-multiplicative involution, and the epsilon-unit bounds. Likewise, def-extended-delta-inclusion says v\_n=I\_n tensor v is a d-inclusion; unpacking its d-homomorphism and two-sided norm clauses gives ||v\_n(T)||<=(1+d)||T||, ||v\_n(ST)-v\_n(S)v\_n(T)||<=d||S||||T||, ||v\_n(I)-I||<=d, and dagger preservation. By the registered external GT-kitaev-standard-diagonal choose D=sum\_j A\_j tensor B\_j with sum\_j||A\_j||||B\_j||=1, XD=DX, and sum\_j A\_jB\_j=I. Put a\_j=I\_n tensor A\_j, b\_j=I\_n tensor B\_j, g\_n(X,Y)=v\_n(XY)-v\_n(X)v\_n(Y), w\_prime\_n(X)=sum\_j v\_n(a\_j)g\_n(b\_j,X), and F\_z(X,Y)=v\_n(X)z(Y)-z(XY)+z(X)v\_n(Y). Direct expansion, using associativity in the source algebra B, gives v\_n(X)g\_n(Y,Z)-g\_n(XY,Z)+g\_n(X,YZ)-g\_n(X,Y)v\_n(Z)=(v\_n(X)v\_n(Y))v\_n(Z)-v\_n(X)(v\_n(Y)v\_n(Z)); the registered epsilon-associativity inequality bounds the right side by epsilon(1+d)\textasciicircum{}3||X||||Y||||Z||<=2epsilon||X||||Y||||Z|| when d+epsilon<=1/4. Expand F\_(w\_prime\_n). Reassociate v\_n(X)(v\_n(a\_j)g\_n(b\_j,Y)), replace v\_n(X)v\_n(a\_j) by v\_n(Xa\_j), and use the amplified identity XD=DX to change the summed expression exactly to sum\_j v\_n(a\_j)g\_n(b\_jX,Y). Reassociate the third term to put v\_n(a\_j) outside and apply the displayed cocycle identity to (b\_j,X,Y), giving sum\_j v\_n(a\_j)(v\_n(b\_j)g\_n(X,Y)) up to its residual. Reassociate again, replace v\_n(a\_j)v\_n(b\_j) by v\_n(a\_jb\_j), use sum\_j a\_jb\_j=I, and finally replace v\_n(I)g\_n(X,Y) by g\_n(X,Y). The seven errors, in this order, are bounded by epsilon(1+d)\textasciicircum{}2*d, (1+epsilon)d\textasciicircum{}2, epsilon(1+d)\textasciicircum{}2*d, 2(1+epsilon)(1+d)epsilon, epsilon(1+d)\textasciicircum{}2*d, (1+epsilon)d\textasciicircum{}2, and ((1+epsilon)d+epsilon)d, times ||X||||Y|| sum\_j||A\_j||||B\_j||. Each bound follows respectively from the explicitly cited epsilon-associativity, d-multiplicativity, product, and epsilon-unit/d-unit inequalities. Since the weight sum is 1 and d+epsilon<=1/4, their sum is at most 100(d\textasciicircum{}2+epsilon)||X||||Y||. Hence ||F\_(w\_prime\_n)(X,Y)-g\_n(X,Y)||<=100(d\textasciicircum{}2+epsilon)||X||||Y|| uniformly in n. \steptag{validated} \\[4pt]
\step{1.3.2} Because v preserves the involution, g\_n(Y\textasciicircum{}dagger,X\textasciicircum{}dagger)\textasciicircum{}dagger=g\_n(X,Y); the definition w\_doubleprime(X)=w\_prime(X\textasciicircum{}dagger)\textasciicircum{}dagger therefore gives ||F\_(w\_doubleprime)-g\_n||<=100(d\textasciicircum{}2+epsilon). Linearity of F gives the same estimate for w. The exact bilinear expansion G\_(v+w)=g-F\_w-w(X)w(Y), together with ||w\_n||<=2d and the product bound, gives ||G\_(v\textasciicircum{}+)\_n(X,Y)||<=105(d\textasciicircum{}2+epsilon)||X||||Y||. Thus C\_m=105 is universal. \steptag{validated} \\[4pt]
\step{1.4} Let v\textasciicircum{}+=v+w, where v is dagger-preserving, w\_doubleprime(X)=w\_prime(X\textasciicircum{}dagger)\textasciicircum{}dagger, and w=(w\_prime+w\_doubleprime)/2. Assume the two estimates supplied separately by the correction branches: sup\_n ||I\_n tensor w|| <= 2d and, at every amplification, ||v\_n\textasciicircum{}+(XY)-v\_n\textasciicircum{}+(X)v\_n\textasciicircum{}+(Y)|| <= 105(d\textasciicircum{}2+epsilon)||X||||Y||. Then v\textasciicircum{}+ is dagger-preserving. Put q=d\textasciicircum{}2+epsilon, mu=105q, x=v\textasciicircum{}+(I\_B), and y=x-I\_A. The level-one displacement estimate and the d-unit bound for v give ||y|| <= 3d, while the multiplicative estimate at X=Y=I\_B gives ||x\textasciicircum{}2-x|| <= mu. Expanding x=I\_A+y and isolating y gives ||y|| <= mu+epsilon(1+epsilon)+2epsilon||y||+(1+epsilon)||y||\textasciicircum{}2. If d+epsilon <= 1/16, then ||y|| <= 3/16 and 2epsilon+(1+epsilon)||y|| <= 83/256; hence ||y|| <= (272/173)(mu+epsilon) < 4(mu+epsilon) <= 424q. Finally I\_n tensor y is the block diagonal direct sum of n copies of y, so the operator-space direct-sum axiom gives ||I\_n tensor y||=||y||. Thus every amplification has unit defect at most 424(d\textasciicircum{}2+epsilon). \steptag{validated} \\[4pt]
\step{1.5} Put q=d\textasciicircum{}2+epsilon and mu\_star=424q. Suppose u:B->A is a dagger-preserving linear map with u\_n=I\_n tensor u such that, for every n, ||u\_n-v\_n||<=2d, ||u\_n(XY)-u\_n(X)u\_n(Y)||<=105q||X||||Y||, and ||u\_n(I)-I||<=mu\_star. Then, if d+epsilon<=e\_step:=min\{1/16,1/(8*424)\}, u is an extended d\textasciicircum{}+-inclusion for d\textasciicircum{}+:=max\{mu\_star,4(mu\_star+epsilon),2mu\_star+epsilon\}<=1700q. Indeed, let M\_n=||u\_n|| and a\_n=inf\_(||X||=1)||u\_n(X)||. Dagger preservation, the target C*-inequality, multiplicativity, and ||X\textasciicircum{}dagger X||=||X||\textasciicircum{}2 give (1-epsilon)M\_n\textasciicircum{}2<=M\_n+mu\_star. For R=1+4(mu\_star+epsilon), (1-epsilon)R\textasciicircum{}2-R-mu\_star=3(mu\_star+epsilon)+(1-epsilon)[4(mu\_star+epsilon)]\textasciicircum{}2-8epsilon(mu\_star+epsilon)>=0 because epsilon<=1/4; hence M\_n<=R. The lower norm bound for v and ||u\_n-v\_n||<=2d give a\_n>=1-3d>=1/2. For unit X, the lower bound a\_n applied to X\textasciicircum{}dagger X, followed by multiplicativity and the target product inequality, gives a\_n<=(1+epsilon)||u\_n(X)||\textasciicircum{}2+mu\_star; taking the infimum yields a\_n<=(1+epsilon)a\_n\textasciicircum{}2+mu\_star. If a\_n>=1 there is nothing to prove, while if a\_n<1, division by a\_n>=1/2 gives 1-a\_n<=epsilon*a\_n+mu\_star/a\_n<=epsilon+2mu\_star. Thus every amplification has the two-sided (1+-d\textasciicircum{}+) norm bounds; its multiplicative defect is at most 105q<=mu\_star<=d\textasciicircum{}+, its unit defect is at most mu\_star<=d\textasciicircum{}+, and it preserves the dagger. Therefore u is an extended d\textasciicircum{}+-inclusion. Since epsilon<=q, d\textasciicircum{}+<=max\{424,4(424+1),2*424+1\}q=1700q. \steptag{validated} \\[4pt]
\end{proofbox}

\end{document}
